\documentclass[11pt]{article}
\usepackage[a4paper,margin=1in]{geometry}
\usepackage[T1]{fontenc}
\usepackage[utf8]{inputenc}
\usepackage{setspace}
\usepackage{hyperref}

\setstretch{1.1}
\setlength{\parskip}{0.6em}
\setlength{\parindent}{0pt}

\begin{document}

\begin{flushright}
June 24, 2026
\end{flushright}

Editors \\
\textit{IEEE Access}

Dear Editors,

We are resubmitting our manuscript (Access-2026-23347) entitled

\textbf{``When Agreement Fails: Visual Anchoring and Multimodal Pseudo-Label
Reliability in Classroom Behavior Recognition''}

following the reviewers' constructive feedback. We thank the Associate Editor
and both reviewers for their careful reading. We have addressed all seven
comments from Reviewer~2 in full; Reviewer~1 raised no technical concerns.
A detailed point-by-point response accompanies this submission.

In addition to the reviewer-driven revisions (restructured Abstract, new
Pipeline Overview, formal visual-anchoring definition, merged tables, and
language polish throughout), we have corrected the first author's affiliation
order: Yan Ma's primary affiliation is now listed as School of Foreign Studies,
Changsha University of Science and Technology (her primary employment), with
School of Education, Hunan Agricultural University as her secondary affiliation
(her PhD enrollment). The revised manuscript also incorporates two new
experiments completed during the revision period:

\begin{enumerate}
  \item \textbf{Ten-model cross-model validation} (expanded from five models
  in the original submission), spanning three model families (Qwen, LLaVA,
  Gemma) across 7B--35B parameter scales. The expanded analysis reveals a
  clear capability boundary on intent-dependent categories and identifies
  Qwen3.5-27B as the most balanced annotator on the hardest dataset.
  \item \textbf{Quality-threshold experiment} using Qwen3.5-27B pseudo-labels
  on TeacherBehavior, which locates the critical annotation accuracy range
  at which pseudo-label fine-tuning begins to exceed the zero-shot baseline---a
  finding with direct practical implications for MLLM-assisted annotation
  pipelines.
\end{enumerate}

The revised manuscript makes four contributions. First, it provides a
reproducible audit protocol linking annotation quality, downstream CLIP
adaptation, retention controls, distribution-shift checks, anchor-proxy
analysis, auxiliary baselines, and a ten-model cross-model robustness check.
Second, it demonstrates that pseudo-label usefulness is conditional on
category-level visual anchoring. Third, it shows that agreement filtering is
not a portable denoising rule. Fourth, it identifies a practical annotation
quality threshold and shows that the selective-routing gain scales with
annotator quality.

We believe the revised manuscript is now substantially stronger and well
suited to \textit{IEEE Access}. All data are publicly available; the code
repository (\url{https://github.com/zhanglizhuo/llm-annotation}) includes
environment specifications, reproduction scripts, and tracked result files.

This submission is original, has not been published previously, and is not
under consideration elsewhere. All authors have approved the revised version
and agree to its resubmission to \textit{IEEE Access}. The authors declare no
competing interests.

Funding statement: This research was funded by the Education Department of
Hunan Province, grant number HNJG-2022-0608.

Data and code availability: The SCB dataset is publicly available at
\url{https://huggingface.co/datasets/wintonYF/SCB-Dataset}. The experiment
code and results are available at
\url{https://github.com/zhanglizhuo/llm-annotation}.

Ethics and privacy: This study is a secondary analysis of a publicly released
benchmark. The manuscript reports aggregate results, makes no re-identification
attempts, and uses privacy-preserving pixelated representative crops where
visual examples are needed.

Thank you for your time and consideration.

Sincerely,

Yan Ma and Lizhuo Zhang

Corresponding author:
Lizhuo Zhang \\
School of Information and Intelligent Science and Technology, \\
Hunan Agricultural University, Changsha 410128, China \\
Email: zhanglizhuo@hunau.edu.cn

\end{document}